\section{Extended Data Tables}\label{app:tables}

This appendix collects extended data behind the six numbered tables of the main text
(Tables~\ref{tab:litchecks}--\ref{tab:floorcensus}) and discusses the three figures that
live only in the appendix (\ref{fig:ak3_plateau}, \ref{fig:hard_ties}, \ref{fig:rl_gap}).
Every value below is read from the same digest files that produced the main-text tables and
figures, and is independently re-derivable from the raw record streams (\ref{app:repro}).

\subsection{Extended per-arm statistics}

Table~\ref{tab:arms} reports median nodes and median path length, over solved presentations
only, for the naive $z=w$ stabilization arms. The full min/max/mean/median breakdown, over
the same solved-only convention, is:

\begin{center}
\small
\begin{tabular}{@{}lll@{}}
\toprule
Arm & Nodes (min / max / mean / median) & Path length (min / max / mean / median) \\
\midrule
baseline & 3 / 77,385 / 1,662.0 / 11 & 2 / 710 / 32.6 / 8 \\
$r_1$    & 5 / 239,578 / 2,802.6 / 13 & 4 / 714 / 33.8 / 12 \\
$r_2$    & 4 / 291,479 / 3,637.1 / 13 & 3 / 715 / 29.9 / 11 \\
$x$      & 4 / 326,326 / 14,680.7 / 10 & 3 / 102 / 12.1 / 9 \\
$y$      & 4 / 491,437 / 21,832.7 / 10 & 3 / 102 / 11.5 / 9 \\
\bottomrule
\end{tabular}
\end{center}

The $x$ and $y$ arms' wide gap between median ($10$) and mean ($14{,}680.7$ / $21{,}832.7$)
node counts, set against a comparatively tight path-length distribution (max $102$), is the
alias-orbit effect noted in \ref{sec:results}: defining $z$ by an already-present generator
does not lengthen solutions, it multiplies equivalent detours the search must exhaust
before terminating.

\subsection{The 97-word literature-grounded bank}

Table~\ref{tab:wordbank} gives family sizes. The full bank --- 95 words across seven
families, plus the two per-target control words $z=r_1,r_2$ added at run time to reach 97 ---
grouped by family in escalation-priority order, is:

\begin{itemize}
  \item \textbf{relhalf} (17): \texttt{xyx, yxy, xxx, yyyy, yxx, xxy, xyy, yyx, XYX, YXY,
    XXX, YYYY, xyxY, Yxyx, xx, yy, yyy}
  \item \textbf{wk} (17): \texttt{yyyyyyyyXyxy, yyyyyyyXyxy, yyyyyyXyxy, yyyyyXyxy,
    yyyyXyxy, yyyXyxy, yyXyxy, yXyxy, Xyxy, YXyxy, YYXyxy, YYYXyxy, YYYYXyxy, YYYYYXyxy,
    YYYYYYXyxy, YYYYYYYXyxy, YYYYYYYYXyxy}
  \item \textbf{wstar} (5): \texttt{YxyX, XyxY, xYXy, yXYx, xYxY}
  \item \textbf{conj} (14): \texttt{x, xyX, yxY, y, Xyx, Yxy, xyxYX, yxyXY, xYxyX, YxyXy,
    xxyXX, yyxYY, XYxyx, XyxYx}
  \item \textbf{comm} (6): \texttt{XYxy, YXyx, xyXY, yxYX, XyXy, xyXYX}
  \item \textbf{ms} (3): \texttt{Xyxyy, yxyx, xyxyy}
  \item \textbf{brute} (33): \texttt{X, Y, xy, xY, XX, Xy, XY, yx, yX, Yx, YX, YY, xxY, xYx,
    xYX, xYY, XXy, XXY, XyX, Xyy, XYx, XYY, yXX, yXy, yXY, yyX, Yxx, YxY, YXX, YXy, YYx, YYX,
    YYY}
\end{itemize}

Each name is the induced word $w$ in $x^{\pm1},y^{\pm1}$ under the encoding of
\ref{app:glossary} (capital letter = inverse); each induces a $z$-relator $z\cdot w^{-1}$,
deduplicated by that induced relator and dropped if it violates the per-relator cap $L=24$.

\subsection{Campaign facet breakdown}

Lane D's $16{,}870$ solve attempts (Table~\ref{tab:lanes}) are drawn from four resumable
facets, every one contributing $0$ solved: D1 ($5{,}866$ solve attempts), D2 ($5{,}350$),
D3 ($5{,}497$), and a resolve pass at the raised per-relator cap $L=40$ ($157$ solve
attempts, probing whether the $L=24$ cap specifically --- as opposed to the total-length cap
--- was pruning a viable descent among the shortest candidates). The four facets sum to the
reported $16{,}870$; $0$ of $16{,}870$ solve attempts succeed across all four.

\subsection{Floor census detail}

The two canonical (signed-relabel) representatives of the length-13 floor of AK(3)'s
AC-class (Table~\ref{tab:floorcensus}) have relators, under the encoding $x\!\to\!1,\
X\!\to\!-1,\ y\!\to\!2,\ Y\!\to\!-2$:

\begin{itemize}
  \item $F$ --- a 2-generator presentation AC-equivalent to AK(3): \texttt{YYxyXX} /
    \texttt{YYYXXyx}, i.e. $y^{-2}xyx^{-2}$ and $y^{-3}x^{-2}yx$ --- the dominant attractor,
    $712/1006$ ($70.8\%$).
  \item AK(3)'s own reduced form: \texttt{YXYxyx} / \texttt{YYYxxxx} --- $294/1006$
    ($29.2\%$).
\end{itemize}

Each census candidate is assigned by canonicalizing its terminal state and checking
equality with one of these two representatives up to signed relabeling (\ref{app:glossary});
no third representative appears anywhere in the 1,006-candidate census.

\subsection{Grid-probe list}

Beyond Lane D's harvested-quotient attempts, the campaign ran $14$ additional grid probes
(Table~\ref{tab:lanes}, Lanes A--C):

\begin{center}
\small
\begin{tabular}{@{}l p{0.36\textwidth} l l@{}}
\toprule
Lane & Probe & Nodes & Floor \\
\midrule
A (MITM) & from AK(3), against 1,177 targets & 2,000,000 & 13 \\
A (MITM) & from P25, against 1,177 targets & 2,000,000 & 13 \\
B (StableSolver) & hero-8 bank, $g\le3$, gen-penalty 2, from P25 & 800,000 & 13 \\
B (StableSolver) & hero-8 bank, $g\le3$, gen-penalty 2, from AK(3) & 800,000 & 13 \\
B (StableSolver) & hero-8 bank, $g\le3$, gen-penalty 1, from AK(3) & 800,000 & 13 \\
B (StableSolver) & hero-8 bank, $g\le4$, gen-penalty 2, from AK(3) & 800,000 & 13 \\
B (StableSolver) & full 95-word bank, $g\le3$, from AK(3) & 300,000 & 13 \\
B (StableSolver) & full 95-word bank, $g\le3$, from P25 & 300,000 & 13 \\
C (trivial-$z$) & rep form, $n=3$ & 1,500,000 & 14 \\
C (trivial-$z$) & textbook form, $n=3$ & 2,000,000 & 14 \\
C (trivial-$z$) & rep form, $n=4$ & 1,000,000 & 15 \\
C (trivial-$z$) & textbook form, $n=4$ & 1,500,000 & 15 \\
C (trivial-$z$) & rep form, $n=5$ & 800,000 & 16 \\
C (trivial-$z$) & textbook form, $n=5$ & 1,000,000 & 16 \\
\bottomrule
\end{tabular}
\end{center}

All 14 probes floor within their searched budgets; $0$ of 14 solve.

\subsection{Discussion of the appendix figures}

\begin{figure}[tbp]
\centering
\includegraphics[width=\textwidth]{ak3_plateau}
\caption{Minimum-total-length reached by the $97$-word bank on both AK(3) forms (trivial
target $=3$); $0$ of $388$ solve attempts succeed. At $10^6$ nodes the distribution
concentrates on the floor --- representative form $\{13\!:\!88,14\!:\!2,15\!:\!7\}$, textbook
form $\{13\!:\!86,14\!:\!2,15\!:\!9\}$ --- a floor of $13$ with a small $14$--$15$ tail that
even the provably-isolatable $w_k$ theory words do not escape.}
\label{fig:ak3_plateau}
\end{figure}

\textbf{Figure~\ref{fig:ak3_plateau}.} At $10^5$ nodes the representative form's floor
histogram is already concentrated ($79$ words at floor $13$, $11$ at $14$, $7$ at $15$);
escalating every unsolved word tenfold to $10^6$ nodes tightens the concentration further
($88/2/7$) rather than dislodging any word past the floor, and the textbook form shows the
identical pattern ($77/11/9$ at $10^5$, sharpening to $86/2/9$ at $10^6$). Every one of the
$97\times2\times2=388$ solve attempts across both forms and both budgets records $0$ solved,
within the searched budgets, against a trivial target of total length $3$; more nodes
concentrate the distribution on the floor rather than escaping it, evidence that the
obstruction is structural rather than a matter of search depth.

\begin{figure}[tbp]
\centering
\includegraphics[width=\textwidth]{hard_ties}
\caption{Word-bank protocol on two hard-but-solvable MS targets. Index $625$: $4$ of $98$
words solve --- controls $z=r_1$ ($77{,}395$ nodes) and $z=r_2$ ($80{,}111$) plus two
relator-derived \texttt{relhalf} words tying those exact node counts. Index $610$: $2$ of
$99$ solve ($z=r_1$ and its \texttt{relhalf} twin, both $61{,}082$). All ${\sim}180$
generic-family runs fail (floor $16$--$25$); no clever word beats or even differs from the
target's own relators.}
\label{fig:hard_ties}
\end{figure}

\textbf{Figure~\ref{fig:hard_ties}.} On the two hard-but-solvable controls, the floor
histograms sit far above AK(3)'s own (index $625$: range $17$--$25$, peaking at $19$ with
$41$ of $98$ words; index $610$: range $16$--$23$, peaking at $19$ with $34$ of $99$ words),
confirming these are genuinely different, easier landscapes with their own internal spread.
Yet the only words that solve are, in both cases, exactly the target's own relators and
their relator-derived \texttt{relhalf} twins, and they solve at identical node counts (index
$625$: $r_1$ and its \texttt{relhalf} twin both at $77{,}395$ nodes; $r_2$ and its twin both
at $80{,}111$; index $610$: $r_1$ and its twin both at $61{,}082$) --- the same search in
disguise, not a shortcut. Within the searched $10^5$-node screen (the specified $10^6$-node
full tier was not run on these generic-family words), no structurally distinct word ever
solves either target.

\begin{figure}[tbp]
\centering
\includegraphics[width=\textwidth]{rl_gap}
\caption{Zero-shot out-of-distribution transfer of a pretrained $2$-generator PPO+beam
policy. It solves $155/155$ of its own training-distribution items (mean path $6.1$, max
$16$) but $0$ of $155$ floor states at beam width $512$; escalation to beam width $2048$
with temperature annealing still solves $0$ of the hardest $30$. The gap is an OOD gap, not
an RL ceiling.}
\label{fig:rl_gap}
\end{figure}

\textbf{Figure~\ref{fig:rl_gap}.} The pretrained policy's own training-distribution path
lengths are short and tightly clustered (mean $6.1$, range $2$--$16$, mode at length $4$:
$31$ of $155$ items), and it solves all $155$ of $155$ solve attempts on that distribution.
Run zero-shot on the $155$ floor and stabilized states at beam width $512$, $0$ of $155$
solve attempts succeed; widening to beam width $2048$ with temperature annealing on the
hardest $30$-state core, $0$ of $30$ solve attempts succeed. Within these two beam-search
configurations, the size of this gap --- perfect in-distribution, zero out-of-distribution,
on states sharing the same 2-generator alphabet the policy was trained on --- is the sharpest
illustration in the campaign that the floor resists every method family tried, learned or
hand-designed, within the budgets searched.
